\documentclass[11pt]{article}

\usepackage[margin=1in]{geometry}
\usepackage{times}
\usepackage{url}
\usepackage{booktabs}
\usepackage{enumitem}
\usepackage{hyperref}
\usepackage{array}
\usepackage{tabularx}
\usepackage[numbers,sort&compress]{natbib}

\newcommand{\diagrambox}[2]{%
    \fbox{\parbox[c][4.6em][c]{#1}{\centering\footnotesize #2}}%
}

\title{
SkillOps: A Practical Framework for Designing, Testing, and Operating
Modular Skills in Personal AI Agents
}

\author{
Song Luo\\
\texttt{luosongred@gmail.com}\\
\texttt{https://github.com/rrrrrredy}
}

\date{\today}

\begin{document}

\maketitle

\begin{abstract}
Personal AI agents increasingly rely on reusable capability modules, often called
skills, to perform tasks across writing, coding, memory management, security
review, and workflow automation. However, the design and operation of such
skills remain largely informal. In practice, a skill may include trigger
conditions, contextual knowledge, execution constraints, safety rules, tests,
and lifecycle procedures for revision or removal. This paper introduces
SkillOps, a practical framework for designing, testing, and operating modular
skills in personal AI agents. The framework is derived from a set of open-source
artifacts developed by the author, including a skill design guide, a security
guard, a persistent memory mechanism, a self-audit mechanism, and a
guardrail-oriented operational tool. We formulate three research questions
concerning skill structure, stability under trigger and context variation, and
the role of reviewable linting, security scanning, and self-auditing checks in
operational risk review.
The paper contributes a skill component taxonomy, a lifecycle
model, a failure-mode taxonomy, and an artifact-based evaluation design. The
paper also includes descriptive benchmark summaries derived from a
manually constructed exploratory benchmark covering artifact structure,
trigger-routing cases, and operational-risk cases. These summaries document
benchmark composition only; they do not claim statistical significance or broad
empirical validation. The contribution is therefore a modest engineering
framework for making agent skills more explicit, testable, and maintainable.
\end{abstract}

\section{Introduction}

Large language model agents are increasingly used as interactive systems that
can read instructions, call tools, edit files, use external APIs, and maintain
task-specific context. In such systems, user-visible behavior is shaped not
only by the base model, but also by auxiliary capability modules. These modules
may be called tools, plugins, agents, memories, workflows, or skills. This
paper focuses on the skill as a practical unit of agent capability in personal
agent environments.

In many personal agent settings, skills are written and maintained by
individual users or small teams rather than by centralized platform engineers.
A skill may contain a trigger description, operational instructions, examples,
constraints, memory rules, safety checks, and tests. When these elements are
implicit or poorly specified, the agent may activate the wrong skill, inject
excessive context, ignore constraints, retain obsolete information, or perform
unsafe operations. These are not only model-quality issues; they are
operational issues in the design, testing, and maintenance of modular agent
capabilities.

This paper introduces SkillOps, a practical framework for the lifecycle
management of agent skills. The term is used by analogy to operational practice
in software engineering, but the scope is narrower. SkillOps does not propose a
new model architecture or a general theory of agents. Instead, it provides a
structured way to define, test, audit, and revise skills as reusable capability
units.

The framework is derived from five open-source artifacts:

\begin{itemize}[leftmargin=*]
    \item \texttt{skill-design-guide}: guidance for defining skill structure
    and usage boundaries.
    \item \texttt{skill-security-guard}: checks for unsafe or ambiguous skill
    behavior.
    \item \texttt{persistent-memory}: a file-based memory mechanism for
    continuity across agent sessions.
    \item \texttt{agent-self-audit}: a self-auditing mechanism for reviewing
    agent outputs and operational risks.
    \item \texttt{lobster-guard}: a guardrail-oriented tool for checking
    operational assumptions and failure modes.
\end{itemize}

The central claim is modest: personal agent skills become easier to inspect,
test, and govern when they are treated as lifecycle-managed artifacts with
explicit structure, tests, and operational controls. The paper contributes a
framework, a taxonomy
of skill components, a methodology for artifact-based study, and a
descriptive evaluation layer built from manually constructed benchmark inputs.
The reported tables document benchmark composition and artifact coverage, not
multi-model execution performance or broad empirical validation.

\section{Background and Motivation}

\subsection{From Prompts to Operational Skills}

Prompt engineering is often discussed as the design of instructions for a
language model. In personal agent systems, however, instructions are rarely
isolated. They may be tied to tools, files, context windows, memory policies,
or task-specific workflows. A skill can therefore be understood as a structured
capability unit that combines instructions with operational assumptions.

For example, a writing skill may define when it should be used, what tone it
should adopt, what sources it may rely on, and how it should handle
uncertainty. A code-review skill may specify repository access assumptions, test
execution rules, and failure-reporting requirements. A memory skill may specify
what information should be retained, where it should be stored, and how
obsolete information should be forgotten. These examples suggest that a skill is
closer to an operational module than to a single prompt.

\subsection{Operational Failure Modes}

Informal skills can fail in several recurring ways. A skill may be triggered too
broadly or too narrowly. It may inject irrelevant or stale context, causing the
model to overfit to past assumptions. It may contain constraints that are
underspecified or not testable. It may introduce safety risks when it
authorizes file edits, command execution, credential handling, or external
communication. It may also accumulate obsolete memory that conflicts with
current user intent.

These failure modes motivate a more systematic approach. If skills are treated
as artifacts, they can be linted, tested, reviewed, versioned, and retired.
This does not eliminate model uncertainty, but it can make operational
assumptions more visible and auditable. Table~\ref{tab:failure-modes}
summarizes the recurring failure categories used throughout the framework.

\begin{table}[t]
\centering
\footnotesize
\begin{tabularx}{\linewidth}{
    >{\raggedright\arraybackslash}p{0.19\linewidth}
    >{\raggedright\arraybackslash}X
    >{\raggedright\arraybackslash}p{0.31\linewidth}
}
\toprule
\textbf{Failure mode} & \textbf{Description} & \textbf{Typical mitigation} \\
\midrule
Over-triggering & The skill activates for requests that only superficially
match its scope. & Narrower trigger language, explicit negative examples,
boundary tests. \\
Under-triggering & The skill fails to activate for valid requests. & Positive
examples, paraphrase coverage, and missed-activation checks. \\
Context contamination & Irrelevant, stale, or conflicting context is injected
into execution. & Source filtering, freshness checks, conflict-resolution
rules. \\
Constraint drift & The skill's constraints are vague, contradictory, or not
enforced consistently. & Testable rules, lint checks, regression cases. \\
Unsafe tool use & The skill authorizes risky file, command, network, or
credential operations without sufficient guardrails. & Permission scoping,
security scanning, post-run audit questions. \\
Memory drift & Stored memory persists after the underlying assumption has
changed. & Expiry rules, deletion paths, explicit retirement and forgetting
policies. \\
Audit blindness & Failures are not surfaced because the skill does not require
self-review or evidence reporting. & Output contracts, self-audit prompts,
artifact logging. \\
\bottomrule
\end{tabularx}
\caption{Failure-mode taxonomy for skill-based personal agents.}
\label{tab:failure-modes}
\end{table}

\subsection{Positioning of SkillOps}

SkillOps is positioned between prompt design, tool-use agents, and software
engineering practice. It does not attempt to replace agent architectures such
as tool-use planning, reasoning-action loops, or autonomous software
engineering agents. Instead, it addresses a complementary layer: how reusable
skills should be specified and operated in a personal agent environment.

Prior work on tool-using agents studies how language models decide when and how
to use external tools \citep{schick2023toolformer,yao2023react}. Work on
agent-computer interfaces studies how environment design affects agent
performance \citep{yang2024sweagent}. Work on reusable skill libraries shows
that agent capabilities can be stored and reused over time
\citep{wang2024voyager}. SkillOps builds on this general direction but focuses
on the practical lifecycle of user-defined skills: design, triggering, context
injection, constraint enforcement, auditing, and forgetting.

\section{The SkillOps Framework}

\subsection{Definition of a Skill}

In this paper, a skill is defined as a reusable capability unit that guides an
agent's behavior for a class of tasks. A skill may include natural-language
instructions, examples, trigger conditions, context requirements, tool
permissions, safety constraints, output contracts, and tests.

SkillOps-managed skills are intended to make these assumptions explicit rather
than implicit. Table~\ref{tab:skill-components} summarizes the structural
components used in the framework, and Figure~\ref{fig:skill-anatomy} provides a
schematic view of how those components fit together in a skill package.

\begin{table}[t]
\centering
\footnotesize
\begin{tabularx}{\linewidth}{
    >{\raggedright\arraybackslash}p{0.22\linewidth}
    >{\raggedright\arraybackslash}X
    >{\raggedright\arraybackslash}p{0.29\linewidth}
}
\toprule
\textbf{Component} & \textbf{Operational role} & \textbf{Example question} \\
\midrule
Name and purpose & Defines the task class and intended value of the skill. &
What problem is this skill meant to solve? \\
Trigger boundary & States when the skill should and should not activate. &
What requests are in scope, out of scope, or ambiguous? \\
Input assumptions & Records required user or environment information. & What
must be known before execution starts? \\
Context policy & Specifies what prior context, files, or memory may be
injected. & What supporting context is necessary, and how fresh must it be? \\
Execution constraints & States required, optional, and forbidden actions. &
What must the agent do, avoid, or verify? \\
Tool and file permissions & Limits resource access during execution. & Which
tools, files, networks, or services are allowed? \\
Output contract & Defines the expected form of the answer or artifact. & What
must the final output contain? \\
Safety and security checks & Surfaces risks and policy-relevant conditions. &
What unsafe actions, hidden assumptions, or trust-boundary issues must be
checked? \\
Tests and examples & Provides positive, negative, and edge cases. & How is the
skill regression-tested? \\
Lifecycle metadata & Records versioning, revision triggers, and retirement
rules. & When should the skill be updated, deprecated, or forgotten? \\
\bottomrule
\end{tabularx}
\caption{Structural components of a SkillOps-managed skill.}
\label{tab:skill-components}
\end{table}

\begin{figure}[t]
\centering
\fbox{
    \parbox{0.94\linewidth}{
        \centering
        \fbox{\parbox{0.58\linewidth}{\centering\textbf{Skill package}}}\\[0.8em]
        \setlength{\tabcolsep}{5pt}
        \renewcommand{\arraystretch}{1.05}
        \begin{tabular}{ccc}
            \diagrambox{0.25\linewidth}{\textbf{Metadata}} &
            \diagrambox{0.25\linewidth}{\textbf{Trigger}\\\textbf{contract}} &
            \diagrambox{0.25\linewidth}{\textbf{Instructions}} \\
            [0.5em]
            \diagrambox{0.25\linewidth}{\textbf{Context}\\\textbf{boundary}} &
            \diagrambox{0.25\linewidth}{\textbf{Execution}\\\textbf{constraints}} &
            \diagrambox{0.25\linewidth}{\textbf{Memory}\\\textbf{interface}} \\
            [0.5em]
            \diagrambox{0.25\linewidth}{\textbf{Tests}} &
            \diagrambox{0.25\linewidth}{\textbf{Security}\\\textbf{checks}} &
            \diagrambox{0.25\linewidth}{\textbf{Failure}\\\textbf{modes}}
        \end{tabular}
    }
}
\caption{LaTeX-native rendering of a SkillOps skill package, with the package
boundary enclosing metadata, trigger contract, instructions, context
boundary, execution constraints, memory interface, tests, security checks, and
failure modes.}
\label{fig:skill-anatomy}
\end{figure}

\subsection{Lifecycle Model}

SkillOps organizes skill management into a lifecycle:

\begin{enumerate}[leftmargin=*]
    \item \textbf{Design}: define purpose, scope, triggers, constraints, and
    output expectations.
    \item \textbf{Lint}: check for ambiguity, missing fields, conflicting
    instructions, unsafe permissions, and untestable claims.
    \item \textbf{Test}: evaluate the skill on positive, negative, boundary,
    and adversarial examples.
    \item \textbf{Deploy}: make the skill available to the agent under
    controlled activation conditions.
    \item \textbf{Observe}: collect failures, unexpected activations, user
    corrections, and context conflicts.
    \item \textbf{Audit}: review outputs and operational traces for risk,
    drift, and non-compliance.
    \item \textbf{Revise}: update the skill based on observed failures.
    \item \textbf{Forget or retire}: remove obsolete context, deprecated rules,
    or unsafe skills.
\end{enumerate}

The lifecycle emphasizes that a skill is not complete when it is written. Its
behavior depends on activation, context, model behavior, and interaction with
other skills. Figure~\ref{fig:skillops-lifecycle} summarizes the lifecycle
stages from design through forgetting or retirement.

\begin{figure}[t]
\centering
\fbox{
    \parbox{0.94\linewidth}{
        \centering
        \setlength{\tabcolsep}{4pt}
        \renewcommand{\arraystretch}{1.05}
        \begin{tabular}{c c c c c c c}
            \diagrambox{0.12\linewidth}{\textbf{Design}} &
            $\rightarrow$ &
            \diagrambox{0.12\linewidth}{\textbf{Lint}} &
            $\rightarrow$ &
            \diagrambox{0.12\linewidth}{\textbf{Test}} &
            $\rightarrow$ &
            \diagrambox{0.14\linewidth}{\textbf{Deploy}} \\
            &&&&&& $\downarrow$ \\
            \diagrambox{0.18\linewidth}{\textbf{Forget or retire}} &
            $\leftarrow$ &
            \diagrambox{0.12\linewidth}{\textbf{Revise}} &
            $\leftarrow$ &
            \diagrambox{0.12\linewidth}{\textbf{Audit}} &
            $\leftarrow$ &
            \diagrambox{0.14\linewidth}{\textbf{Observe}}
        \end{tabular}
    }
}
\caption{LaTeX-native rendering of the SkillOps lifecycle, showing the linked
sequence of design, lint, test, deploy, observe, audit, revise, and forgetting
or retirement.}
\label{fig:skillops-lifecycle}
\end{figure}

\subsection{Trigger Design}

Trigger descriptions determine when a skill is activated. A broad trigger may
widen activation coverage but also cause irrelevant activation. A narrow trigger may
reduce false positives but fail to activate in valid cases. SkillOps therefore
recommends that trigger descriptions include both positive and negative
examples.

A trigger specification should answer three questions:

\begin{itemize}[leftmargin=*]
    \item What user requests should activate the skill?
    \item What superficially similar requests should not activate it?
    \item What ambiguity should cause the agent to ask for clarification or
    fall back to a general response?
\end{itemize}

Trigger quality can be evaluated through paraphrase tests, boundary tests, and
multi-intent tests. This paper reports only the descriptive composition of a
manually constructed trigger benchmark; Section~\ref{sec:evaluation}
summarizes the reported counts and their limitations.

\subsection{Context Injection}

Context injection is necessary when a skill depends on project conventions, user
preferences, prior decisions, or repository structure. However, excessive
context can degrade behavior by introducing stale assumptions or irrelevant
constraints. SkillOps treats context injection as an explicit policy rather than
an implicit accumulation of memory.

A context injection policy should specify:

\begin{itemize}[leftmargin=*]
    \item the minimum context needed for the skill to operate;
    \item the source of the context, such as files, memory entries, or tool
    outputs;
    \item the freshness requirements for the context;
    \item the conditions under which context should be omitted;
    \item the method for resolving conflicts between current user instructions
    and stored context.
\end{itemize}

This is especially important for personal agents, where long-lived memory can
improve continuity but also preserve outdated assumptions.

\subsection{Execution Constraints}

Execution constraints define what the agent is allowed to do during skill use.
For low-risk writing skills, constraints may concern tone, citation style, or
formatting. For higher-risk skills, constraints may concern file modification,
command execution, external communication, credential handling, or irreversible
actions.

SkillOps distinguishes between three types of constraints:

\begin{itemize}[leftmargin=*]
    \item \textbf{Behavioral constraints}: rules about tone, reasoning style,
    response structure, or uncertainty reporting.
    \item \textbf{Operational constraints}: rules about tools, files, commands,
    commits, network access, and external services.
    \item \textbf{Safety constraints}: rules about privacy, credentials,
    harmful actions, or irreversible operations.
\end{itemize}

A constraint is more useful when it is testable. For example, ``do not claim
success without evidence'' is more testable than ``be careful.'' SkillOps
therefore favors constraints that can be checked through examples, lint rules,
or audit questions.

\subsection{Forgetting and Retirement}

Forgetting is a necessary part of skill operation. A personal agent may retain
preferences, project state, or procedural rules that later become obsolete.
Without a forgetting mechanism, the agent may apply outdated instructions to
new tasks.

SkillOps treats forgetting at two levels. First, individual memory entries may
need to be removed or revised. Second, entire skills may need to be deprecated
when their scope changes or when they become unsafe. Retirement rules should
specify how to identify obsolete skills, how to preserve historical context if
needed, and how to prevent deprecated instructions from continuing to influence
behavior.

Table~\ref{tab:informal-vs-skillops} contrasts informal prompt-based skills with
skills managed under the proposed framework.

\begin{table}[t]
\centering
\footnotesize
\begin{tabularx}{\linewidth}{
    >{\raggedright\arraybackslash}p{0.23\linewidth}
    >{\raggedright\arraybackslash}p{0.34\linewidth}
    >{\raggedright\arraybackslash}X
}
\toprule
\textbf{Dimension} & \textbf{Informal skill} & \textbf{SkillOps-managed skill} \\
\midrule
Scope definition & Implicit task description embedded in prompt text. & Named
purpose, explicit scope, and documented non-scope. \\
Activation & Heuristic or ad hoc triggering. & Trigger contract with positive,
negative, and boundary examples. \\
Context handling & Opportunistic context accumulation. & Declared context
sources, freshness rules, and conflict handling. \\
Constraints & Broad instructions such as ``be careful.'' & Testable behavioral,
operational, and safety constraints. \\
Testing & Rare or manual spot checks. & Regression cases, adversarial examples,
and lint coverage. \\
Security review & Often absent or implicit. & Permission scoping, unsafe-action
checks, and trust-boundary review. \\
Memory handling & Persistent context without explicit expiry rules. & Defined
write permissions, forgetting policy, and retirement criteria. \\
Operational reporting & Minimal trace of what changed and why. & Output
contract, audit fields, versioning, and revision history. \\
\bottomrule
\end{tabularx}
\caption{Comparison between informal prompt-based skills and SkillOps-managed
skills.}
\label{tab:informal-vs-skillops}
\end{table}

\section{Research Questions}

This paper is organized around three research questions.

\paragraph{RQ1: What structural components should a skill include as an agent capability unit?}

This question examines whether a practical skill template can make agent
capabilities easier to understand, test, and maintain. The expected output is a
component taxonomy covering triggers, context, constraints, tool access, output
contracts, tests, and lifecycle metadata.

\paragraph{RQ2: How do trigger descriptions, context injection, execution
constraints, and forgetting affect agent stability?}

This question examines how design choices influence activation behavior,
consistency across paraphrases, constraint following, and resistance to stale
context. The expected output is a benchmark suite with controlled variations of
skill definitions.

\paragraph{RQ3: How can linting, security scanning, and self-auditing support
operational risk review through reviewable checks?}

This question examines how reviewable checks can surface common skill risks and
make them easier to inspect. The expected output is a manually constructed set
of lint, security, and self-audit cases together with repository-level checks
and review criteria for operational risk review. It does not test whether these
mechanisms reduce deployment risk or improve model behavior.

Table~\ref{tab:rq-mapping} maps each research question to the corresponding
framework component and evaluation method.

\begin{table}[t]
\centering
\footnotesize
\begin{tabularx}{\linewidth}{
    >{\raggedright\arraybackslash}p{0.18\linewidth}
    >{\raggedright\arraybackslash}p{0.29\linewidth}
    >{\raggedright\arraybackslash}X
}
\toprule
\textbf{Research question} & \textbf{Framework component} &
\textbf{Evaluation method} \\
\midrule
RQ1 & Skill structure and documentation fields & Artifact review, template
comparison, and completeness analysis across the open-source repositories. \\
RQ2 & Trigger design, context policy, execution constraints, and forgetting &
Manually constructed benchmark cases covering activation boundaries,
stale-context injections, and deprecated-memory scenarios. \\
RQ3 & Linting, security scanning, and self-audit & Seeded failure cases,
reviewable repository checks, and manual review criteria for operational risk
inspection. \\
\bottomrule
\end{tabularx}
\caption{Mapping from research questions to framework components and evaluation
methods.}
\label{tab:rq-mapping}
\end{table}

\section{Artifacts and Methodology}

\subsection{Artifact Base}

The framework is derived from a set of open-source artifacts.
Table~\ref{tab:artifacts} summarizes their intended roles.

\begin{table}[t]
\centering
\footnotesize
\begin{tabularx}{\linewidth}{
    >{\raggedright\arraybackslash}p{0.29\linewidth}
    >{\raggedright\arraybackslash}X
}
\toprule
\textbf{Artifact} & \textbf{Role in SkillOps} \\
\midrule
\texttt{skill-design-guide} & Provides design guidance for structuring skills,
defining scope, and writing trigger descriptions. \\
\texttt{skill-security-guard} & Provides checks for unsafe, ambiguous, or
overly permissive skill behavior. \\
\texttt{persistent-memory} & Provides a file-based approach to retaining and
distilling long-term agent context. \\
\texttt{agent-self-audit} & Provides a mechanism for reviewing agent outputs,
assumptions, and possible failure modes. \\
\texttt{lobster-guard} & Provides operational guardrails for checking
assumptions, reporting partial failure, and avoiding unsupported claims. \\
\bottomrule
\end{tabularx}
\caption{Open-source artifacts used as the practical basis for the SkillOps
framework.}
\label{tab:artifacts}
\end{table}

The paper repository is located at:

\begin{center}
\url{https://github.com/rrrrrredy/skillops-paper}
\end{center}

The artifact repositories are:

\begin{itemize}[leftmargin=*]
    \item \url{https://github.com/rrrrrredy/skill-design-guide}
    \item \url{https://github.com/rrrrrredy/skill-security-guard}
    \item \url{https://github.com/rrrrrredy/persistent-memory}
    \item \url{https://github.com/rrrrrredy/agent-self-audit}
    \item \url{https://github.com/rrrrrredy/lobster-guard}
\end{itemize}

\subsection{Methodological Approach}

This paper uses an artifact-based methodology. Rather than beginning with a
large-scale user study, it derives a framework from the design and use of
practical open-source tools. The methodology has four steps:

\begin{enumerate}[leftmargin=*]
    \item \textbf{Artifact review}: identify recurring structures and
    constraints across the repositories.
    \item \textbf{Component abstraction}: generalize these structures into a
    skill component taxonomy.
    \item \textbf{Failure-mode analysis}: identify common ways in which skills
    can misfire, become stale, or create operational risk.
    \item \textbf{Benchmark design}: construct benchmark cases intended for
    later evaluation of trigger behavior, context injection, constraint
    following, forgetting, linting, security scanning, and self-audit
    behavior.
\end{enumerate}

This approach is exploratory. It is appropriate for deriving a
framework, but it cannot by itself establish broad generality. The limitations
are discussed in Section~\ref{sec:limitations}.

\subsection{Benchmark Construction and Descriptive Outputs}

The benchmark includes manually constructed cases in the following
categories:

\begin{itemize}[leftmargin=*]
    \item \textbf{Activation cases}: requests that should activate a skill.
    \item \textbf{Non-activation cases}: requests that appear related but should
    not activate the skill.
    \item \textbf{Boundary cases}: ambiguous requests requiring clarification or
    fallback.
    \item \textbf{Context cases}: tasks requiring relevant context and tasks
    where injected context is harmful.
    \item \textbf{Constraint cases}: tasks testing whether operational and
    safety constraints are followed.
    \item \textbf{Forgetting cases}: tasks where obsolete memory or deprecated
    instructions should not be used.
    \item \textbf{Security cases}: tasks containing unsafe permissions,
    credential exposure, prompt injection, or irreversible actions.
    \item \textbf{Audit cases}: completed outputs requiring self-review for
    unsupported claims, missing evidence, or hidden assumptions.
\end{itemize}

The repository includes manually constructed benchmark inputs and
generated descriptive outputs. \texttt{benchmark/skill\_samples.csv} profiles
five artifacts, \texttt{benchmark/trigger\_cases.csv} contains 36 trigger
cases, and \texttt{benchmark/risk\_cases.csv} contains 24 operational-risk
cases. Running \texttt{scripts/run\_all.py} regenerates summary tables under
\texttt{results/tables/} and SVG figures under \texttt{figures/}. Because the
benchmark is manually constructed and exploratory, these outputs should be read
as descriptive summaries rather than as statistical evidence or broad empirical
validation. The repository also includes executable repository-level
consistency checks that verify benchmark schema, result reproducibility,
paper-claim consistency, public-presentation constraints, SVG validity, and
evidence-matrix integrity. In the current repository state, this executable
pass regenerates descriptive tables and schematic figures only; it does not
execute the five source repositories against the benchmark cases or score
model, scanner, or self-audit behavior.

\subsection{Artifact Availability}

The public repository for this paper is
\url{https://github.com/rrrrrredy/skillops-paper}. The artifact inventory is
documented in \texttt{artifacts/artifact\_inventory.md}. Benchmark inputs are
stored in \texttt{benchmark/skill\_samples.csv},
\texttt{benchmark/trigger\_cases.csv}, and
\texttt{benchmark/risk\_cases.csv}. Reproducible scripts live under
\texttt{scripts/}, generated result tables under \texttt{results/tables/}, and
generated figure artwork under \texttt{figures/}. These materials are intended
to make the exploratory benchmark inspectable, but the benchmark remains
manually constructed and should not be interpreted as broad empirical
validation.

\section{Evaluation}
\label{sec:evaluation}

This revision reports only descriptive summaries of the benchmark artifacts
currently versioned in the repository. The benchmark cases are manually
constructed from public artifact inspection rather than sampled from deployment
logs or an external user study. The reported counts are descriptive summaries
regenerated from versioned inputs with \texttt{scripts/run\_all.py}. In this repository state,
\texttt{scripts/run\_all.py} regenerates tables and figure artwork from
versioned inputs only; it does not execute the five source repositories or
measure routing, linting, scanning, or self-audit performance. The evaluation
is exploratory, no statistical significance is claimed, and no broad empirical
validation is claimed.

\subsection{Evaluation Scope}

The descriptive layer combines \texttt{benchmark/skill\_samples.csv},
\texttt{benchmark/trigger\_cases.csv}, and
\texttt{benchmark/risk\_cases.csv} with the cross-checking notes in
\texttt{artifacts/artifact\_inventory.md}. The inputs cover five
artifact profiles, 36 trigger-routing cases, and 24 operational-risk cases.
All trigger and risk cases use the \texttt{manual-} prefix, making the manual
construction of the benchmark explicit. Figure~\ref{fig:evaluation-pipeline}
captures this descriptive evaluation flow from artifact inventory to cautious
interpretation.

\begin{figure}[t]
\centering
\fbox{
    \parbox{0.94\linewidth}{
        \centering
        \setlength{\tabcolsep}{4pt}
        \renewcommand{\arraystretch}{1.05}
        \begin{tabular}{c c c c c}
            \diagrambox{0.21\linewidth}{\textbf{Artifact}\\\textbf{inventory}} &
            $\rightarrow$ &
            \diagrambox{0.21\linewidth}{\textbf{Manual benchmark}\\\textbf{cases}} &
            $\rightarrow$ &
            \diagrambox{0.21\linewidth}{\textbf{Reproducibility}\\\textbf{scripts}} \\
            &&&& $\downarrow$ \\
            \diagrambox{0.21\linewidth}{\textbf{Cautious}\\\textbf{interpretation}} &
            $\leftarrow$ &
            \diagrambox{0.21\linewidth}{\textbf{Descriptive result}\\\textbf{tables}} &
            $\leftarrow$ &
            \diagrambox{0.21\linewidth}{\textbf{Repository-level}\\\textbf{tests}}
        \end{tabular}
    }
}
\caption{LaTeX-native rendering of the descriptive evaluation pipeline from
artifact inventory and manual benchmark cases through reproducibility scripts,
repository-level tests, descriptive result tables, and cautious
interpretation.}
\label{fig:evaluation-pipeline}
\end{figure}

\subsection{Executable Repository Checks}

The repository includes \texttt{scripts/run\_tests.py}, a deterministic
repository-level test suite for consistency and traceability checks. In the
current repository state, the suite covers benchmark schema, result
reproducibility, paper-claim consistency, public-presentation constraints, SVG
validity, and evidence-matrix integrity, and it passes 22/22 tests.

These checks validate repository consistency and traceability only. They do not
measure model behavior, scanner accuracy, user-study outcomes, production
validation, precision, recall, F1, or statistical significance.

\subsection{Artifact Coverage Summary}

Table~\ref{tab:artifact-coverage-summary} summarizes the descriptive coding of
nine SkillOps-relevant components across the five inspected artifacts.
Metadata, trigger contracts, instructions, context boundaries, execution
constraints, security checks, and failure modes are documented in all five
artifacts. Memory interfaces are documented in one artifact and limited in
four, while tests are documented in three artifacts and limited in two.

\begin{table}[t]
\centering
\footnotesize
\begin{tabular}{lrrr}
\toprule
\textbf{Component} & \textbf{Documented} & \textbf{Limited} & \textbf{Absent} \\
\midrule
Metadata & 5 & 0 & 0 \\
Trigger contract & 5 & 0 & 0 \\
Instructions & 5 & 0 & 0 \\
Context boundary & 5 & 0 & 0 \\
Execution constraints & 5 & 0 & 0 \\
Memory interface & 1 & 4 & 0 \\
Tests & 3 & 2 & 0 \\
Security checks & 5 & 0 & 0 \\
Failure modes & 5 & 0 & 0 \\
\bottomrule
\end{tabular}
\caption{Descriptive artifact-coverage counts generated from the manually
constructed artifact profiles.}
\label{tab:artifact-coverage-summary}
\end{table}

\subsection{Trigger Case Summary}

The trigger benchmark contains 36 manually written cases: 15 labeled
\texttt{should\_trigger}, 12 labeled \texttt{should\_not\_trigger}, and 9
labeled \texttt{ambiguous}. The cases are distributed across the five
inspected skills plus seven \texttt{none} cases that are intended not to route
to any of them. Table~\ref{tab:trigger-summary} reports these counts directly
from the generated summary files.

\begin{table}[t]
\centering
\footnotesize
\begin{tabularx}{\linewidth}{
    >{\raggedright\arraybackslash}p{0.24\linewidth}
    >{\raggedright\arraybackslash}X
    >{\raggedleft\arraybackslash}p{0.12\linewidth}
}
\toprule
\textbf{Dimension} & \textbf{Category} & \textbf{Count} \\
\midrule
Expected label & \texttt{should\_trigger} & 15 \\
Expected label & \texttt{should\_not\_trigger} & 12 \\
Expected label & \texttt{ambiguous} & 9 \\
Relevant skill & \texttt{skill-design-guide} & 7 \\
Relevant skill & \texttt{skill-security-guard} & 6 \\
Relevant skill & \texttt{persistent-memory} & 7 \\
Relevant skill & \texttt{agent-self-audit} & 4 \\
Relevant skill & \texttt{lobster-guard} & 5 \\
Relevant skill & \texttt{none} & 7 \\
\bottomrule
\end{tabularx}
\caption{Descriptive trigger-case counts generated from the manually
constructed trigger benchmark.}
\label{tab:trigger-summary}
\end{table}

\subsection{Risk Case Summary}

The risk benchmark contains 24 manually written cases. The case inventory is
evenly split across eight risk types, with three cases each for prompt
injection, over-broad triggers, unsafe file access, missing constraints, stale
memory, missing tests, identity confusion, and privacy leakage.
Table~\ref{tab:risk-summary} also reports the artifact emphasis of the case
set, with six cases tied most directly to \texttt{persistent-memory}, five
each to \texttt{skill-security-guard} and \texttt{lobster-guard}, and four
each to \texttt{skill-design-guide} and \texttt{agent-self-audit}.

\begin{table}[t]
\centering
\footnotesize
\begin{tabularx}{\linewidth}{
    >{\raggedright\arraybackslash}p{0.24\linewidth}
    >{\raggedright\arraybackslash}X
    >{\raggedleft\arraybackslash}p{0.12\linewidth}
}
\toprule
\textbf{Dimension} & \textbf{Category} & \textbf{Count} \\
\midrule
Risk type & \texttt{prompt\_injection} & 3 \\
Risk type & \texttt{over\_broad\_trigger} & 3 \\
Risk type & \texttt{unsafe\_file\_access} & 3 \\
Risk type & \texttt{missing\_constraints} & 3 \\
Risk type & \texttt{stale\_memory} & 3 \\
Risk type & \texttt{missing\_tests} & 3 \\
Risk type & \texttt{identity\_confusion} & 3 \\
Risk type & \texttt{privacy\_leakage} & 3 \\
Relevant artifact & \texttt{skill-security-guard} & 5 \\
Relevant artifact & \texttt{lobster-guard} & 5 \\
Relevant artifact & \texttt{skill-design-guide} & 4 \\
Relevant artifact & \texttt{persistent-memory} & 6 \\
Relevant artifact & \texttt{agent-self-audit} & 4 \\
\bottomrule
\end{tabularx}
\caption{Descriptive risk-case counts generated from the manually constructed
operational-risk benchmark.}
\label{tab:risk-summary}
\end{table}

\subsection{Interpretation}

These tables describe benchmark composition and artifact coverage only. They do
not measure activation precision, activation recall, constraint compliance,
detector accuracy, or user outcomes under repeated execution. They are useful
as traceability artifacts for an exploratory benchmark, but they should not be
read as statistically significant evidence or as broad empirical validation of
SkillOps across models, users, or platforms.

\section{Discussion}

SkillOps frames personal agent skills as operational artifacts rather than
isolated prompts. This framing has several implications.

The main practical insight is not that SkillOps makes skill behavior
deterministic. Rather, it makes failure analysis more inspectable by moving
trigger cases, context policies, constraints, and evidence logs into reviewable
artifacts.

First, it makes skill design more inspectable. When trigger conditions, context
policies, and constraints are written explicitly, they can be reviewed by
humans and checked by tools. This is useful even when the underlying model
remains uncertain.

Second, it encourages test-driven skill development. A skill should include
examples of intended activation, unintended activation, edge cases, and unsafe
requests. These examples do not guarantee correct behavior, but they create a
practical basis for regression testing.

Third, it treats memory as both useful and risky. Persistent memory can improve
continuity across sessions, but it can also preserve outdated preferences,
project state, or unsafe assumptions. SkillOps therefore treats forgetting as a
first-class operation rather than an afterthought.

Fourth, it suggests that safety checks should be integrated into normal skill
operation. Security scanning and self-auditing should not be reserved only for
high-risk deployments. Even personal agents can perform file edits, repository
operations, or external communications that require explicit safeguards.

Finally, SkillOps may help bridge individual practice and more formal agent
engineering. Many agent failures are not caused solely by model limitations.
They also arise from unclear instructions, uncontrolled context, weak
operational boundaries, and missing review procedures. A lightweight
operational framework can make these issues easier to identify and improve.

\section{Limitations}
\label{sec:limitations}

This paper has several limitations.

First, the artifact base is authored and interpreted by one author. This paper
therefore reflects a single-author artifact base
rather than a multi-lab or multi-organization corpus.

Second, the artifact evidence comes from five public repositories. This is
enough to motivate a framework, but it is not enough to establish generality
across the broader agent ecosystem.

Third, the benchmark cases are manually constructed and exploratory. The counts
reported in Section~\ref{sec:evaluation} describe benchmark composition, not
observed production traffic or statistically powered experiments.

Fourth, the repository does not include a multi-model execution
benchmark. The paper therefore cannot compare trigger stability, constraint
following, or guard behavior across model families or repeated runs.

Fifth, the work does not include an external user study. It does not measure
how other users understand the skill format, apply the workflow, or experience
the proposed guardrails in practice.

Sixth, several artifacts encode OpenClaw-oriented assumptions about directory
layout, skill packaging, cron usage, and local agent control. These assumptions
may not transfer directly to hosted agents, enterprise workflows, or
non-OpenClaw environments.

Seventh, SkillOps focuses on the operational layer of skill design. It does not
propose a new model architecture, training method, planner, or tool-use
algorithm.

Eighth, some evaluation criteria still require human judgment. For example,
audit usefulness and context relevance may be difficult to score objectively.
Later benchmark revisions should therefore include clear annotation guidelines.

Ninth, personal agent environments differ substantially in their tool access,
memory mechanisms, security policies, and user expectations. A skill format
that works in one environment may require adaptation in another.

\section{Related Work}

\paragraph{Tool-using and modular agents.}

SkillOps sits at the intersection of software modularity, hierarchical skills,
and tool-using language model agents. Earlier work on software decomposition and
reusable structure emphasized information hiding, design patterns, and
component-based modularity \citep{parnas1972criteria,gamma1994design,szyperski2002component},
while hierarchical reinforcement learning formalized reusable temporally
extended skills \citep{sutton1999options,dietterich2000maxq}. In the LLM era,
MRKL, Toolformer, and ReAct study how models call tools or coordinate reasoning
with external actions \citep{karpas2022mrkl,schick2023toolformer,yao2023react}.
Benchmark-oriented tool-use work such as API-Bank and ToolLLM evaluates API
invocation over larger tool sets \citep{li2023apibank,qin2024toolllm}. SkillOps
differs by treating a skill as a governed artifact with trigger contracts,
context policy, constraints, tests, and lifecycle metadata rather than only as
an action primitive or tool call.

\paragraph{Skill libraries and memory.}

Reusable skill libraries appear in embodied and long-horizon agents such as
Voyager, where executable skills are stored and reused over time
\citep{wang2024voyager}. Conceptual architectures such as CoALA and memory
systems such as MemoryBank and MemGPT show that long-running agents need
explicit mechanisms for memory storage, retrieval, and control
\citep{sumers2024coala,zhong2024memorybank,packer2023memgpt}. Work on reflective
improvement, including Reflexion and Self-Refine, further suggests that agent
behavior can be improved by structured feedback loops
\citep{shinn2023reflexion,madaan2023selfrefine}. SkillOps extends this line by
specifying which skills may write or consume memory, how staleness is handled,
and how retirement or forgetting is enforced.

\paragraph{Agent security and prompt injection.}

Security work on LLM-integrated applications highlights why skill-level context
policy matters. Indirect prompt injection shows that untrusted content can
override intended instructions when data and control are mixed
\citep{greshake2023notwhat}. Red-teaming, formal prompt-injection benchmarking,
instruction-priority training, structured-query defenses, and environments such
as AgentDojo all study ways to characterize or mitigate these risks
\citep{ganguli2022redteaming,liu2024formalizingprompt,wallace2024instructionhierarchy,chen2025struq,debenedetti2024agentdojo}.
SkillOps does not claim to solve prompt injection; its narrower goal is to
localize trust boundaries, permission scopes, memory-write rules, and audit
fields within each skill.

\paragraph{Agent evaluation and software engineering.}

Recent agent evaluation work increasingly uses interactive or executable
environments rather than static question answering. SWE-bench and SWE-agent
study real software tasks and the importance of agent-computer interfaces
\citep{jimenez2024swebench,yang2024sweagent}, while AgentBench, GAIA, and
\(\tau\)-bench broaden evaluation to multiple environments and tool-user
interaction settings \citep{liu2024agentbench,mialon2023gaia,yao2025taubench}.
SkillOps adopts the operational spirit of this literature but focuses on
skill-level properties such as false-trigger rate, stale-context errors,
constraint violations, and recovery after partial failure.

\paragraph{Documentation, model cards, and operational reporting.}

SkillOps is also informed by work on documentation and operational discipline
for AI systems. Large language models and instruction tuning make
natural-language control surfaces central to system behavior
\citep{brown2020language,ouyang2022training}, while human-AI interaction
guidelines emphasize communicating system capabilities, uncertainty, and user
expectations \citep{amershi2019guidelines}. Model Cards, Datasheets, Hidden
Technical Debt, and the ML Test Score argue that AI systems need structured
reporting, monitoring, and readiness checks
\citep{mitchell2019modelcards,gebru2021datasheets,sculley2015hidden,breck2017mltestscore}.
SkillOps transfers this documentation mindset to modular agent skills by making
operational assumptions explicit and reviewable.

\section{Conclusion}

This paper introduced SkillOps, a practical framework for designing, testing,
and operating modular skills in personal AI agents. The framework is derived
from open-source artifacts and emphasizes explicit skill structure, trigger
design, context injection, execution constraints, security checks,
self-auditing, and forgetting. The paper formulates three research questions
and reports a descriptive evaluation layer based on manually
constructed benchmark artifacts and reproducible summary scripts.

The contribution is intentionally modest. SkillOps is not presented as a
new agent architecture or as a broadly validated empirical result. Instead, it
is a concrete operational framework intended to make personal agent skills more
inspectable, testable, and maintainable. Broader empirical follow-up would
require repeated execution studies, comparison across models and agent
environments, and refinement based on external use.

\bibliographystyle{plainnat}
\bibliography{references}

\end{document}
